\label{sec:disc}

\subsection{Class Imbalance}

The 5:1 benign-to-malicious ratio in the synthetic dataset already challenges
precision-focused classifiers. In practice, the ratio may be 500:1 or higher: an operator
session is minutes long and occurs hundreds of times per day, while a sophisticated
intrusion may span a single multi-hour campaign that generates a handful of malicious
sessions. The \texttt{class\_weight=balanced} setting compensates at the classifier level,
but extremely sparse minority-class examples may still degrade recall. Techniques such as
SMOTE oversampling, cost-sensitive learning, or anomaly-detection framing (training only
on benign sessions and flagging deviations) should be evaluated as the class ratio
approaches operational reality.

\subsection{Concept Drift in Production}

Windows Event Log distributions are non-stationary: software updates introduce new event
IDs (or retire existing ones), organisational restructuring changes user and hostname
populations, and seasonal operational patterns alter inter-event timing. A model trained
on a static snapshot will experience concept drift \cite{sommer2010outside}. Mitigation
strategies include periodic retraining on sliding windows of recent labelled sessions,
online learning with drift detectors (e.g., ADWIN), and calibration against analyst-reviewed
alert outcomes. The session-level feature abstraction provides a degree of stability—event
ID counts are invariant to log format changes—but time-delta features may shift following
infrastructure updates.

\subsection{GDPR and Data Retention Constraints}

Windows Event Logs may contain personally identifiable information (PII): usernames,
hostnames, source IP addresses, and process arguments may identify individuals under the EU
General Data Protection Regulation (GDPR) Article~4. ICS operators face a tension between
the NERC CIP-007 requirement to retain audit logs for up to 90~days and GDPR data
minimisation and storage-limitation principles. Practical mitigations include
pseudonymisation of user and host identifiers at collection time, encryption of stored logs
with role-based key management (consistent with IEC~62351), and automated deletion
workflows aligned with the GDPR retention schedule. The feature engineering pipeline in this
work is inherently privacy-preserving: session-level aggregates (counts, statistics) contain
no individual-level identifiers.

\subsection{Limitations}

The primary limitation of this study is the use of synthetic data. The synthetic generator
captures the structural characteristics of the four attack patterns but does not model
evasion behaviours (e.g., time-throttled reconnaissance, log-clearing attempts). Detection
performance on operational logs with real noise, label uncertainty, and adversarial evasion
is expected to be lower. A secondary limitation is the single-site, single-time-period
evaluation; generalisability across grid operators with different audit policies and host
populations is untested.
